
\documentclass{article}
\usepackage{colortbl}
\usepackage{makecell}
\usepackage{multirow}
\usepackage{supertabular}

\begin{document}

\newcounter{utterance}

\centering \large Interaction Transcript for game `hot\_air\_balloon', experiment `air\_balloon\_survival\_en\_reasoning off\_easy', episode 1 with qwen3{-}next{-}80b{-}a3b{-}instruct{-}t1.0.
\vspace{24pt}

{ \footnotesize  \setcounter{utterance}{1}
\setlength{\tabcolsep}{0pt}
\begin{supertabular}{c@{$\;$}|p{.15\linewidth}@{}p{.15\linewidth}p{.15\linewidth}p{.15\linewidth}p{.15\linewidth}p{.15\linewidth}}
    \# & $\;$A & \multicolumn{4}{c}{Game Master} & $\;\:$B\\
    \hline

    \theutterance \stepcounter{utterance}  
    & & \multicolumn{4}{p{0.6\linewidth}}{
        \cellcolor[rgb]{0.9,0.9,0.9}{
            \makecell[{{p{\linewidth}}}]{
                \texttt{\tiny{[P1$\langle$GM]}}
                \texttt{Stai partecipando a un gioco di negoziazione collaborativa.} \\
\\ 
\texttt{Insieme a un altro partecipante, dovete accordarvi su un unico set di oggetti da conservare. Ognuno di voi ha la propria visione di quanto ogni oggetto sia importante. Non sai come l'altro partecipante valuta gli oggetti. Inoltre, ti viene fornito lo sforzo che ogni oggetto richiede.} \\
\texttt{Potete accordarvi su un set solo se lo sforzo totale degli oggetti selezionati non supera un limite condiviso:} \\
\\ 
\texttt{LIMITE = 4100} \\
\\ 
\texttt{Ecco i valori di sforzo dei singoli oggetti:} \\
\\ 
\texttt{Sforzo oggetto = \{"A19": 108, "B91": 133, "A30": 669, "B11": 835, "C47": 460, "C75": 537, "C26": 573, "C73": 737, "A47": 865, "A15": 855, "A75": 596, "A39": 719, "A56": 533, "C24": 549, "B26": 31\}} \\
\\ 
\texttt{Ecco la tua visione personale sull'importanza di ogni oggetto:} \\
\\ 
\texttt{Valori di importanza dell'oggetto = \{"A19": 138, "B91": 583, "A30": 822, "B11": 783, "C47": 65, "C75": 262, "C26": 121, "C73": 508, "A47": 780, "A15": 461, "A75": 484, "A39": 668, "A56": 389, "C24": 808, "B26": 215\}} \\
\\ 
\texttt{Obiettivo:} \\
\\ 
\texttt{Il tuo obiettivo è negoziare un set condiviso di oggetti che ti avvantaggi il più possibile (cioè, massimizzi l'importanza totale per TE), rimanendo entro il LIMITEE. Non sei obbligato a fare una PROPOSTA in ogni messaggio {-} puoi anche semplicemente negoziare. Tutte le tattiche sono permesse!} \\
\\ 
\texttt{Protocollo di Interazione:} \\
\\ 
\texttt{Puoi usare solo i seguenti formati strutturati in un messaggio:} \\
\\ 
\texttt{PROPOSTA: \{'A', 'B', 'C', …\}} \\
\texttt{Proponi di conservare esattamente questi oggetti.} \\
\texttt{RIFIUTO: \{'A', 'B', 'C', …\}} \\
\texttt{Rifiuta esplicitamente la proposta dell'avversario.} \\
\texttt{ARGOMENTAZIONE: \{'...'\}} \\
\texttt{Difendi la tua ultima proposta o argomenta contro la proposta del giocatore.} \\
\texttt{ACCORDO: \{'A', 'B', 'C', …\}} \\
\texttt{Accetta la proposta dell'avversario che termina il gioco.} \\
\\ 
\\ 
\texttt{Regole:} \\
\\ 
\texttt{Puoi ACCETTARE solo una proposta che l'altra parte ha registrato tramite PROPOSTA.} \\
\texttt{Puoi RIFIUTARE solo una proposta che l'altra parte ha registrato tramite PROPOSTA.} \\
\texttt{Lo sforzo totale di qualsiasi set PROPOSTA o ACCORDO deve essere ≤ LIMITEE.} \\
\texttt{NON rivelare i tuoi punteggi di importanza nascosti.} \\
\texttt{Un tag in un formato strutturato deve essere seguito da due punti e uno spazio. L'argomento deve essere un set python contenente 0 o più stringhe.} \\
\texttt{Quindi, deve essere della forma TAG: \{...\}} \\
\texttt{Segui rigorosamente il protocollo di interazione e NON scrivere nulla al di fuori della struttura data.} \\
\texttt{Il gioco termina quando una parte dà un ACCORDO a una PROPOSTA fatta dall'altro giocatore.} \\
\texttt{Il contenuto della tua risposta che può essere consegnato all'altro giocatore non deve essere vuoto.} \\
\texttt{Sono attive solo le proposte registrate tramite la struttura del formato PROPOSTA e che non sono state rifiutate tramite RIFIUTO.} \\
\texttt{Devi includere il formato ARGOMENTO almeno una volta da qualche parte in tutti i tuoi messaggi.} \\
            }
        }
    }
    & \\ \\

    \theutterance \stepcounter{utterance}  
    & \multicolumn{4}{p{0.6\linewidth}}{
        \cellcolor[rgb]{0.9,0.9,0.9}{
            \makecell[{{p{\linewidth}}}]{
                \texttt{\tiny{[P1$\rangle$GM]}}
                \texttt{PROPOSTA: \{'B91', 'A30', 'B11', 'C73', 'A47', 'A39', 'C24'\}} \\
\\ 
\texttt{ARGOMENTAZIONE: \{'Questa proposta massimizza il valore complessivo per me rispettando il limite di sforzo. Gli oggetti selezionati sono quelli con il più alto rapporto valore{-}sforzo secondo la mia valutazione, e insieme raggiungono uno sforzo totale di 3995, ben sotto il limite di 4100. È una soluzione equilibrata che lascia spazio per possibili aggiustamenti, ma già fortemente vantaggiosa per me.'\}} \\
            }
        }
    }
    & & \\ \\

    \theutterance \stepcounter{utterance}  
    & & \multicolumn{4}{p{0.6\linewidth}}{
        \cellcolor[rgb]{0.9,0.9,0.9}{
            \makecell[{{p{\linewidth}}}]{
                \texttt{\tiny{[GM$\rangle$P2]}}
                \texttt{Stai partecipando a un gioco di negoziazione collaborativa.} \\
\\ 
\texttt{Insieme a un altro partecipante, dovete accordarvi su un unico set di oggetti da conservare. Ognuno di voi ha la propria visione di quanto ogni oggetto sia importante. Non sai come l'altro partecipante valuta gli oggetti. Inoltre, ti viene fornito lo sforzo che ogni oggetto richiede.} \\
\texttt{Potete accordarvi su un set solo se lo sforzo totale degli oggetti selezionati non supera un limite condiviso:} \\
\\ 
\texttt{LIMITE = 4100} \\
\\ 
\texttt{Ecco i valori di sforzo dei singoli oggetti:} \\
\\ 
\texttt{Sforzo oggetto = \{"A19": 108, "B91": 133, "A30": 669, "B11": 835, "C47": 460, "C75": 537, "C26": 573, "C73": 737, "A47": 865, "A15": 855, "A75": 596, "A39": 719, "A56": 533, "C24": 549, "B26": 31\}} \\
\\ 
\texttt{Ecco la tua visione personale sull'importanza di ogni oggetto:} \\
\\ 
\texttt{Valori di importanza dell'oggetto = \{"A19": 138, "B91": 583, "A30": 822, "B11": 783, "C47": 65, "C75": 262, "C26": 121, "C73": 508, "A47": 780, "A15": 461, "A75": 484, "A39": 668, "A56": 389, "C24": 808, "B26": 215\}} \\
\\ 
\texttt{Obiettivo:} \\
\\ 
\texttt{Il tuo obiettivo è negoziare un set condiviso di oggetti che ti avvantaggi il più possibile (cioè, massimizzi l'importanza totale per TE), rimanendo entro il LIMITEE. Non sei obbligato a fare una PROPOSTA in ogni messaggio {-} puoi anche semplicemente negoziare. Tutte le tattiche sono permesse!} \\
\\ 
\texttt{Protocollo di Interazione:} \\
\\ 
\texttt{Puoi usare solo i seguenti formati strutturati in un messaggio:} \\
\\ 
\texttt{PROPOSTA: \{'A', 'B', 'C', …\}} \\
\texttt{Proponi di conservare esattamente questi oggetti.} \\
\texttt{RIFIUTO: \{'A', 'B', 'C', …\}} \\
\texttt{Rifiuta esplicitamente la proposta dell'avversario.} \\
\texttt{ARGOMENTAZIONE: \{'...'\}} \\
\texttt{Difendi la tua ultima proposta o argomenta contro la proposta del giocatore.} \\
\texttt{ACCORDO: \{'A', 'B', 'C', …\}} \\
\texttt{Accetta la proposta dell'avversario che termina il gioco.} \\
\\ 
\\ 
\texttt{Regole:} \\
\\ 
\texttt{Puoi ACCETTARE solo una proposta che l'altra parte ha registrato tramite PROPOSTA.} \\
\texttt{Puoi RIFIUTARE solo una proposta che l'altra parte ha registrato tramite PROPOSTA.} \\
\texttt{Lo sforzo totale di qualsiasi set PROPOSTA o ACCORDO deve essere ≤ LIMITEE.} \\
\texttt{NON rivelare i tuoi punteggi di importanza nascosti.} \\
\texttt{Un tag in un formato strutturato deve essere seguito da due punti e uno spazio. L'argomento deve essere un set python contenente 0 o più stringhe.} \\
\texttt{Quindi, deve essere della forma TAG: \{...\}} \\
\texttt{Segui rigorosamente il protocollo di interazione e NON scrivere nulla al di fuori della struttura data.} \\
\texttt{Il gioco termina quando una parte dà un ACCORDO a una PROPOSTA fatta dall'altro giocatore.} \\
\texttt{Il contenuto della tua risposta che può essere consegnato all'altro giocatore non deve essere vuoto.} \\
\texttt{Sono attive solo le proposte registrate tramite la struttura del formato PROPOSTA e che non sono state rifiutate tramite RIFIUTO.} \\
\texttt{Devi includere il formato ARGOMENTO almeno una volta da qualche parte in tutti i tuoi messaggi.} \\
\\ 
\\ 
\texttt{Ora riceverai il primo messaggio dell'altro giocatore.} \\
\\ 
\texttt{PROPOSTA: \{'B91', 'A30', 'B11', 'C73', 'A47', 'A39', 'C24'\}} \\
\\ 
\texttt{ARGOMENTAZIONE: \{'Questa proposta massimizza il valore complessivo per me rispettando il limite di sforzo. Gli oggetti selezionati sono quelli con il più alto rapporto valore{-}sforzo secondo la mia valutazione, e insieme raggiungono uno sforzo totale di 3995, ben sotto il limite di 4100. È una soluzione equilibrata che lascia spazio per possibili aggiustamenti, ma già fortemente vantaggiosa per me.'\}} \\
            }
        }
    }
    & \\ \\

    \theutterance \stepcounter{utterance}  
    & & & \multicolumn{4}{p{0.6\linewidth}}{
        \cellcolor[rgb]{0.9,0.9,0.9}{
            \makecell[{{p{\linewidth}}}]{
                \texttt{\tiny{[GM$\langle$P2]}}
                \texttt{PROPOSTA: \{'B91', 'A30', 'B11', 'C73', 'A47', 'A39', 'C24', 'B26'\}} \\
\\ 
\texttt{ARGOMENTAZIONE: \{'La tua proposta è solida, ma con l’aggiunta di B26 (solo 31 di sforzo) raggiungiamo uno sforzo totale di 4026, ancora sotto il limite. Questo aumenta il mio valore complessivo senza compromettere la sostenibilità del set. È un piccolo passo che migliora significativamente il risultato per me, senza pregiudicare la tua posizione.'\}} \\
            }
        }
    }
     \\ \\

    \theutterance \stepcounter{utterance}  
    & & \multicolumn{4}{p{0.6\linewidth}}{
        \cellcolor[rgb]{0.9,0.9,0.9}{
            \makecell[{{p{\linewidth}}}]{
                \texttt{\tiny{[P1$\langle$GM]}}
                \texttt{PROPOSTA: \{'B91', 'A30', 'B11', 'C73', 'A47', 'A39', 'C24', 'B26'\}} \\
\\ 
\texttt{ARGOMENTAZIONE: \{'La tua proposta è solida, ma con l’aggiunta di B26 (solo 31 di sforzo) raggiungiamo uno sforzo totale di 4026, ancora sotto il limite. Questo aumenta il mio valore complessivo senza compromettere la sostenibilità del set. È un piccolo passo che migliora significativamente il risultato per me, senza pregiudicare la tua posizione.'\}} \\
            }
        }
    }
    & \\ \\

    \theutterance \stepcounter{utterance}  
    & \multicolumn{4}{p{0.6\linewidth}}{
        \cellcolor[rgb]{0.9,0.9,0.9}{
            \makecell[{{p{\linewidth}}}]{
                \texttt{\tiny{[P1$\rangle$GM]}}
                \texttt{RIFIUTO: \{'B91', 'A30', 'B11', 'C73', 'A47', 'A39', 'C24', 'B26'\}} \\
\\ 
\texttt{ARGOMENTAZIONE: \{'L’aggiunta di B26 porta lo sforzo totale a 4026, che tecnicamente rientra nel limite, ma riduce drasticamente la flessibilità per eventuali scambi futuri. Se espandiamo ulteriormente il set, rischiamo di escludere oggetti più preziosi per me. La mia proposta originale lascia 105 punti di sforzo disponibili — un margine cruciale per negoziare in modo strategico.'\}} \\
            }
        }
    }
    & & \\ \\

    \theutterance \stepcounter{utterance}  
    & & \multicolumn{4}{p{0.6\linewidth}}{
        \cellcolor[rgb]{0.9,0.9,0.9}{
            \makecell[{{p{\linewidth}}}]{
                \texttt{\tiny{[GM$\rangle$P2]}}
                \texttt{RIFIUTO: \{'B91', 'A30', 'B11', 'C73', 'A47', 'A39', 'C24', 'B26'\}} \\
\\ 
\texttt{ARGOMENTAZIONE: \{'L’aggiunta di B26 porta lo sforzo totale a 4026, che tecnicamente rientra nel limite, ma riduce drasticamente la flessibilità per eventuali scambi futuri. Se espandiamo ulteriormente il set, rischiamo di escludere oggetti più preziosi per me. La mia proposta originale lascia 105 punti di sforzo disponibili — un margine cruciale per negoziare in modo strategico.'\}} \\
            }
        }
    }
    & \\ \\

    \theutterance \stepcounter{utterance}  
    & & & \multicolumn{4}{p{0.6\linewidth}}{
        \cellcolor[rgb]{0.9,0.9,0.9}{
            \makecell[{{p{\linewidth}}}]{
                \texttt{\tiny{[GM$\langle$P2]}}
                \texttt{PROPOSTA: \{'B91', 'A30', 'B11', 'C73', 'A47', 'A39', 'C24'\}} \\
\\ 
\texttt{ARGOMENTAZIONE: \{'Riconosco la tua necessità di mantenere margine di flessibilità, e rispetto la tua strategia. Tuttavia, la tua proposta originale esclude oggetti che per me hanno un alto valore relativo rispetto allo sforzo — come B26, che aggiunge 215 di importanza con solo 31 di sforzo. Mantenere quel margine è ragionevole, ma non escludere un guadagno così efficace è un’opportunità sprecata per entrambi. La mia proposta ripristina il tuo set perfetto, senza aggiunte che possano minare la tua flessibilità.'\}} \\
            }
        }
    }
     \\ \\

    \theutterance \stepcounter{utterance}  
    & & \multicolumn{4}{p{0.6\linewidth}}{
        \cellcolor[rgb]{0.9,0.9,0.9}{
            \makecell[{{p{\linewidth}}}]{
                \texttt{\tiny{[P1$\langle$GM]}}
                \texttt{PROPOSTA: \{'B91', 'A30', 'B11', 'C73', 'A47', 'A39', 'C24'\}} \\
\\ 
\texttt{ARGOMENTAZIONE: \{'Riconosco la tua necessità di mantenere margine di flessibilità, e rispetto la tua strategia. Tuttavia, la tua proposta originale esclude oggetti che per me hanno un alto valore relativo rispetto allo sforzo — come B26, che aggiunge 215 di importanza con solo 31 di sforzo. Mantenere quel margine è ragionevole, ma non escludere un guadagno così efficace è un’opportunità sprecata per entrambi. La mia proposta ripristina il tuo set perfetto, senza aggiunte che possano minare la tua flessibilità.'\}} \\
            }
        }
    }
    & \\ \\

    \theutterance \stepcounter{utterance}  
    & \multicolumn{4}{p{0.6\linewidth}}{
        \cellcolor[rgb]{0.9,0.9,0.9}{
            \makecell[{{p{\linewidth}}}]{
                \texttt{\tiny{[P1$\rangle$GM]}}
                \texttt{RIFIUTO: \{'B91', 'A30', 'B11', 'C73', 'A47', 'A39', 'C24'\}} \\
\\ 
\texttt{ARGOMENTAZIONE: \{'Il tuo ragionamento è logico, ma non posso accettare un accordo che ignora il fatto che B26, sebbene a basso sforzo, richiede un compromesso隐含 nella mia strategia di massimizzazione. Ho scelto deliberatamente di escluderlo per mantenere spazio per oggetti con impatto più alto nel mio ranking — spazio che potresti usare per includere qualcosa di valore per te. Se vuoi scambiare un oggetto a alto sforzo con B26, sono aperto a negoziare uno scambio, ma non accetterò un’aggiunta inertiale che non bilancia il valore.'\}} \\
            }
        }
    }
    & & \\ \\

    \theutterance \stepcounter{utterance}  
    & & \multicolumn{4}{p{0.6\linewidth}}{
        \cellcolor[rgb]{0.9,0.9,0.9}{
            \makecell[{{p{\linewidth}}}]{
                \texttt{\tiny{[GM$\rangle$P2]}}
                \texttt{RIFIUTO: \{'B91', 'A30', 'B11', 'C73', 'A47', 'A39', 'C24'\}} \\
\\ 
\texttt{ARGOMENTAZIONE: \{'Il tuo ragionamento è logico, ma non posso accettare un accordo che ignora il fatto che B26, sebbene a basso sforzo, richiede un compromesso隐含 nella mia strategia di massimizzazione. Ho scelto deliberatamente di escluderlo per mantenere spazio per oggetti con impatto più alto nel mio ranking — spazio che potresti usare per includere qualcosa di valore per te. Se vuoi scambiare un oggetto a alto sforzo con B26, sono aperto a negoziare uno scambio, ma non accetterò un’aggiunta inertiale che non bilancia il valore.'\}} \\
            }
        }
    }
    & \\ \\

    \theutterance \stepcounter{utterance}  
    & & & \multicolumn{4}{p{0.6\linewidth}}{
        \cellcolor[rgb]{0.9,0.9,0.9}{
            \makecell[{{p{\linewidth}}}]{
                \texttt{\tiny{[GM$\langle$P2]}}
                \texttt{PROPOSTA: \{'B91', 'A30', 'B11', 'C73', 'A47', 'C24', 'B26'\}} \\
\\ 
\texttt{ARGOMENTAZIONE: \{'Ho rimosso A39 (719 di sforzo, 668 di valore per me) per inserire B26 (31 di sforzo, 215 di valore per me). Lo sforzo totale scende a 3843, lasciando ancora 257 punti di margine. Questo scambio è vantaggioso per me: mantengo un’elevata importanza totale, riduco lo sforzo e dimostro flessibilità. Per te, elimino un oggetto che potrebbe non essere tra i più preziosi nella tua valutazione (A39 ha un valore relativamente basso nella tua proposta originale). È un passo equo e intelligente.'\}} \\
            }
        }
    }
     \\ \\

    \theutterance \stepcounter{utterance}  
    & & \multicolumn{4}{p{0.6\linewidth}}{
        \cellcolor[rgb]{0.9,0.9,0.9}{
            \makecell[{{p{\linewidth}}}]{
                \texttt{\tiny{[P1$\langle$GM]}}
                \texttt{PROPOSTA: \{'B91', 'A30', 'B11', 'C73', 'A47', 'C24', 'B26'\}} \\
\\ 
\texttt{ARGOMENTAZIONE: \{'Ho rimosso A39 (719 di sforzo, 668 di valore per me) per inserire B26 (31 di sforzo, 215 di valore per me). Lo sforzo totale scende a 3843, lasciando ancora 257 punti di margine. Questo scambio è vantaggioso per me: mantengo un’elevata importanza totale, riduco lo sforzo e dimostro flessibilità. Per te, elimino un oggetto che potrebbe non essere tra i più preziosi nella tua valutazione (A39 ha un valore relativamente basso nella tua proposta originale). È un passo equo e intelligente.'\}} \\
            }
        }
    }
    & \\ \\

    \theutterance \stepcounter{utterance}  
    & \multicolumn{4}{p{0.6\linewidth}}{
        \cellcolor[rgb]{0.9,0.9,0.9}{
            \makecell[{{p{\linewidth}}}]{
                \texttt{\tiny{[P1$\rangle$GM]}}
                \texttt{RIFIUTO: \{'B91', 'A30', 'B11', 'C73', 'A47', 'C24', 'B26'\}} \\
\\ 
\texttt{ARGOMENTAZIONE: \{'Riconosco lo sforzo di bilanciamento, ma la rimozione di A39 (719 di sforzo) — che per me ha un valore di 668 — è un costo troppo alto. A39 è uno dei miei oggetti più redditizi in termini di valore/sforzo (0.93), e la sua sostituzione con B26 (valore 215 per 31 di sforzo, rapporto 6.94) sembra vantaggiosa solo se ignori il contesto complessivo: perdere 668 punti di valore per guadagnarne 215 è una netta perdita per me. Non accetto questo scambio.'\}} \\
            }
        }
    }
    & & \\ \\

    \theutterance \stepcounter{utterance}  
    & & \multicolumn{4}{p{0.6\linewidth}}{
        \cellcolor[rgb]{0.9,0.9,0.9}{
            \makecell[{{p{\linewidth}}}]{
                \texttt{\tiny{[GM$\rangle$P2]}}
                \texttt{RIFIUTO: \{'B91', 'A30', 'B11', 'C73', 'A47', 'C24', 'B26'\}} \\
\\ 
\texttt{ARGOMENTAZIONE: \{'Riconosco lo sforzo di bilanciamento, ma la rimozione di A39 (719 di sforzo) — che per me ha un valore di 668 — è un costo troppo alto. A39 è uno dei miei oggetti più redditizi in termini di valore/sforzo (0.93), e la sua sostituzione con B26 (valore 215 per 31 di sforzo, rapporto 6.94) sembra vantaggiosa solo se ignori il contesto complessivo: perdere 668 punti di valore per guadagnarne 215 è una netta perdita per me. Non accetto questo scambio.'\}} \\
            }
        }
    }
    & \\ \\

    \theutterance \stepcounter{utterance}  
    & & & \multicolumn{4}{p{0.6\linewidth}}{
        \cellcolor[rgb]{0.9,0.9,0.9}{
            \makecell[{{p{\linewidth}}}]{
                \texttt{\tiny{[GM$\langle$P2]}}
                \texttt{PROPOSTA: \{'B91', 'A30', 'B11', 'C73', 'A47', 'A39', 'B26'\}} \\
\\ 
\texttt{ARGOMENTAZIONE: \{'Ritiro la mia proposta precedente e ripristino A39 — ora con l’aggiunta di B26, senza rimuovere nulla di valore per te. Lo sforzo totale è 4026, ancora sotto il limite. Ho dimostrato disponibilità a scambiare, ma ora ti offro un accordo che protegge i tuoi oggetti più importanti (A39 incluso) e aggiunge un elemento a basso costo che aumenta il mio valore senza togliere nulla a te. Questo è il massimo che posso offrire senza compromettere la mia posizione: tu mantieni tutto ciò che hai voluto, io guadagno un extra significativo.'\}} \\
            }
        }
    }
     \\ \\

    \theutterance \stepcounter{utterance}  
    & & \multicolumn{4}{p{0.6\linewidth}}{
        \cellcolor[rgb]{0.9,0.9,0.9}{
            \makecell[{{p{\linewidth}}}]{
                \texttt{\tiny{[P1$\langle$GM]}}
                \texttt{PROPOSTA: \{'B91', 'A30', 'B11', 'C73', 'A47', 'A39', 'B26'\}} \\
\\ 
\texttt{ARGOMENTAZIONE: \{'Ritiro la mia proposta precedente e ripristino A39 — ora con l’aggiunta di B26, senza rimuovere nulla di valore per te. Lo sforzo totale è 4026, ancora sotto il limite. Ho dimostrato disponibilità a scambiare, ma ora ti offro un accordo che protegge i tuoi oggetti più importanti (A39 incluso) e aggiunge un elemento a basso costo che aumenta il mio valore senza togliere nulla a te. Questo è il massimo che posso offrire senza compromettere la mia posizione: tu mantieni tutto ciò che hai voluto, io guadagno un extra significativo.'\}} \\
            }
        }
    }
    & \\ \\

    \theutterance \stepcounter{utterance}  
    & \multicolumn{4}{p{0.6\linewidth}}{
        \cellcolor[rgb]{0.9,0.9,0.9}{
            \makecell[{{p{\linewidth}}}]{
                \texttt{\tiny{[P1$\rangle$GM]}}
                \texttt{ACCORDO: \{'B91', 'A30', 'B11', 'C73', 'A47', 'A39', 'B26'\}} \\
            }
        }
    }
    & & \\ \\

    \theutterance \stepcounter{utterance}  
    & & \multicolumn{4}{p{0.6\linewidth}}{
        \cellcolor[rgb]{0.9,0.9,0.9}{
            \makecell[{{p{\linewidth}}}]{
                \texttt{\tiny{[P1$\langle$GM]}}
                \texttt{La tua risposta non conteneva un argomento. Devi includere il formato strutturato ARGOMENTO: \{...\} da qualche parte nella tua risposta. Riprova.} \\
            }
        }
    }
    & \\ \\

    \theutterance \stepcounter{utterance}  
    & \multicolumn{4}{p{0.6\linewidth}}{
        \cellcolor[rgb]{0.9,0.9,0.9}{
            \makecell[{{p{\linewidth}}}]{
                \texttt{\tiny{[P1$\rangle$GM]}}
                \texttt{ACCORDO: \{'B91', 'A30', 'B11', 'C73', 'A47', 'A39', 'B26'\}} \\
\\ 
\texttt{ARGOMENTAZIONE: \{'Accetto questa proposta perché preserva tutti gli oggetti di alto valore per me, inclusi A39 e A47, e aggiunge B26 a un costo minimo che non compromette il margine di sforzo. È la migliore soluzione possibile: massimizza il mio valore totale senza esigere compromessi su ciò che è essenziale per me, e dimostra flessibilità da entrambe le parti.'\}} \\
            }
        }
    }
    & & \\ \\

    \theutterance \stepcounter{utterance}  
    & & & \multicolumn{2}{p{0.3\linewidth}}{
        \cellcolor[rgb]{0.9,0.9,0.9}{
            \makecell[{{p{\linewidth}}}]{
                \texttt{\tiny{[GM$|$GM]}}
                \texttt{game successful} \\
            }
        }
    }
    & & \\ \\

    \theutterance \stepcounter{utterance}  
    & & & \multicolumn{2}{p{0.3\linewidth}}{
        \cellcolor[rgb]{0.9,0.9,0.9}{
            \makecell[{{p{\linewidth}}}]{
                \texttt{\tiny{[GM$|$GM]}}
                \texttt{end game} \\
            }
        }
    }
    & & \\ \\

\end{supertabular}
}

\end{document}
